\subsection{An actual-source normalization from eighteen pairs to four}

\begin{statement}{D23.C30}{Proved}{Cubing reduces the existential mixed gate to four pairs}
\end{statement}
If a normalized actual full-GT element with
$(\rho,a_3,a_2)=(0,1,0)$ exists, then one exists whose
degree-1344 Hurwitz data $(V,R)$ is one of the four pairs
\[
                        (1,0),\quad(19,0),\quad(23,1),\quad(27,2)
                                                               \tag{D23.C30.1}
\]
in the fixed eighteen-pair numbering.  In each of these
four pairs, both $V$ and $U=R^{-1}V$ are involutions.
This is an existential normalization by replacing the
element by its cube; it does not reject the other
fourteen possibilities for an element before normalization.

\begin{proof}
The multiplier-one normalized free action associated to
$h=(1,f_h)$ is
\[
            \varphi_h(x)=x,\qquad
            \varphi_h(y)=f_h^{-1}y f_h.
\]
Its composition law is
\[
              f_{hk}=f_h\varphi_h(f_k).                   \tag{D23.C30.2}
\]
This follows directly by composing the two automorphisms
on $y$; uniqueness of the commutator normalization removes
the possible power of $y$.  It is the multiplier-one
specialization, in the present action convention, of the
standard GT composition formula of Lochak--Schneps,
\emph{Inventiones Mathematicae} \textbf{127} (1997),
571--600, equation (2) on printed page 572.

For the mixed finite pair, actual covariance gives
$VxV^{-1}=RxR^{-1}$ and $VyV^{-1}=y$, where $R=f_h(x,y)$.
Thus $U=R^{-1}V$ implements $\varphi_h$ on the finite
mixed image.  Evaluation of (D23.C30.2), or induction on
powers of $h$, yields
\[
                    R_{h^n}=V^nU^{-n}.                    \tag{D23.C30.3}
\]
The actual Hurwitz action is a group action, so its
permutation for $h^n$ is $V^n$.  These statements apply
to powers of a single actual element and do not assume
that different surviving finite candidates can be
realized simultaneously.

Exact enumeration of the eighteen necessary pairs gives
orders dividing six for both $V$ and $U$.  Applying
(D23.C30.3) at $n=3$ gives exactly the four pairs
(D23.C30.1).  All four have $V^2=U^2=1$.
For example the three pairs above each of base indices
$1,19,23,27$ cube respectively to $(1,0),(19,0),(23,1),(27,2)$;
all six pairs above $53,61$ cube to $(1,0)$.
The certificate also checks every odd power in the
finite cyclic group generated by each pair, and each
remains in the original eighteen-element list.

Finally $\rho\bmod3$, $a_2$ and $a_3$ are characters on
the multiplier-one group.  Cubing preserves
$\rho=0\bmod3$, $a_2=0$, and $a_3=1$.
Consequently $h^3$ is still a normalized mixed element,
and its finite data belongs to (D23.C30.1).  This proves
the smaller existential gate.
\end{proof}

\paragraph{Certificate.}
\path{agents/characters/code/hurwitz_candidate_power_closure.py}
computes all orders, every odd power, and the complete
cube map.  Its JSON/NPZ records have the same stem.
The input eighteen-pair list is the main actual-source
necessary reduction, including the uniqueness of the
first value.  The power normalization uses no new
lifting assertion.
